\documentclass[11pt]{article}
\usepackage[T1]{fontenc}
\usepackage{lmodern,microtype}
\usepackage[margin=1in]{geometry}
\usepackage{amsmath,amssymb,amsthm,mathtools}
\usepackage[hidelinks]{hyperref}
\newtheorem{theorem}{Theorem}
\newtheorem{lemma}[theorem]{Lemma}
\newtheorem{proposition}[theorem]{Proposition}
\newtheorem{corollary}[theorem]{Corollary}
\DeclareMathOperator{\per}{per}
\DeclareMathOperator{\rank}{rank}
\DeclareMathOperator{\Ree}{Re}
\DeclareMathOperator{\Imm}{Im}
\newcommand{\contr}{\operatorname{contr}}
\newcommand{\one}{\mathbf 1}
\newcommand{\CP}{\mathbb {CP}}
\title{Sharp logarithmic growth of permanental cofactor spectra\\and an optimal rank-two endpoint constant}
\author{Yongxian Zhang\thanks{School of Computer Science and Engineering,
South China University of Technology, Guangzhou, China.
Email: \href{mailto:mxymmxym1@gmail.com}{mxymmxym1@gmail.com}.
ORCID: \href{https://orcid.org/0009-0000-3864-3536}{0009-0000-3864-3536}.}}
\date{8 October 2026}
\hypersetup{pdfauthor={Yongxian Zhang}}
\setlength{\emergencystretch}{3em}
\begin{document}
\maketitle
\begin{abstract}
For a Hermitian positive-semidefinite matrix $A$, let
$C(A)_{ij}=a_{ij}\per A(i\mid j)$, with row $i$ and column $j$ deleted.
The largest possible ratio $\lambda_{\max}(C(A))/\per A$ among matrices of
order $N$ is asymptotic to $\log N$. The corresponding maximum over real
Rayleigh directions, $\lambda_{\max}(\Ree C(A))/\per A$, is asymptotic to
$\tfrac12\log N$. Both lower limits already hold within rank-two correlation
matrices at every sufficiently large order. We give an explicit dyadic family,
exact finite Rayleigh certificates, and a finite random-sign construction
attaining the leading constants. We also determine the sharp asymptotic ramp
Rayleigh constant $3\pi^2/16$ and its rank-two $q$-permanent endpoint counterpart
$\pi^2/8$. These are cofactor extremal results; Lieb's character-dominance and
Marcus's permanent-of-permanents conjectures remain unresolved here.
\end{abstract}

\section{Definitions, statements, and attribution}
All Hermitian inner products are linear in their first argument. The minor
$A(i\mid j)$ is not transposed. For a PSD matrix with $\per A>0$, put
\[
 C(A)_{ij}=a_{ij}\per A(i\mid j).
\]
The real part of a matrix is entrywise. Thus $\lambda_{\max}(\Ree C(A))$ is the
maximum of $w^*C(A)w$ over real unit vectors, even when $A$ has non-real entries.
It is not a supremum restricted to real matrices.

Let $R_N,R_N^{\mathbb R}$ be the suprema of the two normalized eigenvalues over
all complex Hermitian PSD matrices of order $N$ with positive permanent.
Let $R_N^{(2)},R_N^{(2,\mathbb R)}$ restrict these suprema to rank-exactly-two
correlation matrices, for $N\ge2$.
\begin{theorem}\label{thm:main}
As $N\to\infty$,
\[
 \frac{R_N}{\log N}\longrightarrow1,\qquad
 \frac{R_N^{\mathbb R}}{\log N}\longrightarrow\frac12,\qquad
 \frac{R_N^{(2)}}{\log N}\longrightarrow1,\qquad
 \frac{R_N^{(2,\mathbb R)}}{\log N}\longrightarrow\frac12.
\]
The first two limits are unchanged if the suprema are restricted to
positive-definite correlation matrices.
\end{theorem}

Drury's order-eight example already disproved the cofactor eigenvalue
conjecture with constant one \cite{drury}; that disproof is not new here.
The first-compound indicator inequality used for our upper bound is a classical
refinement of Lieb's inequality, appearing in Pate's work \cite{pate96}.
We supply its full contraction proof. The connection of cofactor directions
with local Hadamard perturbations is also established in the literature
\cite{pioge}. Our assertions concern dimension-dependent extrema and the
sharp ramp constant. A separate source audit records the finite search and
access limitations; no exhaustive priority claim is made. In particular no
unchecked all-$k$ reading of the abstract of \cite{pate08} is used.

\section{The first-compound indicator inequality and the upper bounds}
For homogeneous polynomials use the Bargmann--Fock inner product
\[
 \left\langle\sum_\alpha p_\alpha x^\alpha,
 \sum_\alpha q_\alpha x^\alpha\right\rangle_B
 =\sum_\alpha\alpha!p_\alpha\overline{q_\alpha}.
\]
If $A$ is the Gram matrix of rows $v_i$ and $\ell_i$ their linear forms, the
pairing expansion gives
\begin{equation}\label{eq:fock}
 \per A=\left\|\prod_i\ell_i\right\|_B^2,
 \qquad C(A)=\operatorname{Gram}
 \left(v_i\otimes\prod_{j\ne i}\ell_j\right)_{i=1}^N.
\end{equation}
Thus $C(A)\succeq0$. The Laplace expansion gives $C(A)\one=(\per A)\one$.
These statements also follow for singular matrices directly from their Gram
factorizations.

\begin{lemma}\label{lem:indicator}
For every subset $S\subseteq\{1,\ldots,N\}$,
\[
 \one_S^*C(A)\one_S\le |S|\per A.
\]
\end{lemma}
\begin{proof}
Write $a=|S|$, $b=N-a$. The empty and full subsets are immediate. Relabel so
$S=\{1,\ldots,a\}$. Let $u$ and $v$ be the symmetrizations of the two row
tensors, of degrees $a$ and $b$, respectively, with symmetrization defined by
averaging over permutations. Classify a permutation by the number $j$ of
indices mapped from $S$ to $S^c$. The permutations in this class form a double
coset of $S_a\times S_b$ with cardinality
\[
 a!b!\binom aj\binom bj.
\]
Indeed choose the $j$ cross rows and columns in each block, match the cross
sets, and match the two remaining sets; the count simplifies to the displayed
number. Averaging the original tensor on both sides over $S_a\times S_b$
replaces it by $u\otimes v$. For the representative swapping the first $j$
coordinates of the two blocks, summing those coordinates shows that its
matrix coefficient at $u\otimes v$ is
$\|u\mathbin{\contr_j}\overline v\|^2\ge0$. Here the contraction pairs $j$
coordinates and leaves the other coordinates in the tensor product of the
underlying space and its conjugate. In coordinates it is the squared norm
of the array obtained by summing the products of $j$ matching entries of $u$
and $\overline v$; this is the usual two-block symmetric-tensor contraction.
Consequently the sum of Gram products in this class is
\[
 t_j=a!b!\binom aj\binom bj
       \|u\mathbin{\contr_j}\overline v\|^2\ge0,
 \qquad \per A=\sum_jt_j.
\]
Expanding $\sum_{i,l\in S}a_{il}\per A(i\mid l)$ counts a permutation once
for each $i\in S$ with its image also in $S$. There are exactly $a-j$ such
indices in the $j$th class. Hence the sum equals
$\sum_j(a-j)t_j\le a\sum_jt_j$.
\end{proof}

\begin{lemma}\label{lem:layer}
If a $d$ by $d$ PSD Hermitian matrix $M$ satisfies
$b^*Mb\le\|b\|^2$ for all binary $b$, then, with
$S_d=\sum_{j=1}^d(\sqrt j-\sqrt{j-1})^2$,
\[
 \lambda_{\max}(M)\le4S_d\le4+H_{d-1},\qquad
 \lambda_{\max}(\Ree M)\le2S_d\le2+\tfrac12H_{d-1}.
\]
Here $H_0=0$.
\end{lemma}
\begin{proof}
Sort a nonnegative vector $x$ so $x_1\ge\cdots\ge x_d\ge0$, and set
$x_{d+1}=0$. In the seminorm $\|M^{1/2}\cdot\|$, the nested-indicator
expansion and summation by parts give
\[
 \|M^{1/2}x\|
 \le\sum_j(x_j-x_{j+1})\sqrt j
 =\sum_jx_j(\sqrt j-\sqrt{j-1})
 \le\sqrt{S_d}\|x\|.
\]
Splitting a real vector into its positive and negative parts bounds its
seminorm by $\sqrt{2S_d}$ times its norm. Splitting a complex vector into the
four positive and negative real and imaginary parts gives $2\sqrt{S_d}$.
These are the two spectral bounds. For $j\ge2$,
$(\sqrt j-\sqrt{j-1})^2\le1/(4(j-1))$, giving the harmonic estimates.
\end{proof}
Applying the lemmas to $C(A)/\per A$ proves
\begin{equation}\label{eq:upper}
 R_N\le4+H_{N-1},\qquad R_N^{\mathbb R}\le2+\tfrac12H_{N-1}.
\end{equation}

\section{A finite ring construction attaining the leading constants}
\begin{lemma}\label{lem:entropy}
For $b>1$ put $E(b)=b^{b/(b-1)}/(b-1)$. If positive numbers
$d_1>\cdots>d_s$ satisfy $d_i/d_{i+1}\ge b$ and have sum $j$, then
\[
 \frac{j^j}{\prod_i d_i^{d_i}}\le E(b)^j.
\]
\end{lemma}
\begin{proof}
Index $p_i=d_i/j$ starting at zero. The first $k$ terms sum to at least
$p_{k-1}(b^k-1)/(b-1)$, whereas the remaining tail is at most
$p_{k-1}/(b-1)$. Thus $\sum_{i\ge k}p_i\le b^{-k}$, and
$\sum_i i p_i\le1/(b-1)$. Nonnegativity of relative entropy with respect to
$g_i=(1-1/b)b^{-i}$ gives
\[
 -\sum_i p_i\log p_i
 \le-\log(1-1/b)+(\log b)/(b-1)=\log E(b).
\]
Multiplication by $j$ and exponentiation finish the proof. Zero probabilities
past the finite support cause no problem in the relative entropy inequality.
\end{proof}

Fix $b>1$, $0<\delta<1$, and $\eta>0$. Choose a positive real number
\[
 C>\frac{E(b)}{e(1-\delta)},\qquad
 \theta=\frac{E(b)}{Ce(1-\delta)}<1,\qquad q=\theta^2.
\]
Choose an integer $d_0\ge2$ so large that
\begin{equation}\label{eq:d0}
 2e\exp\!\left(\frac{q^{d_0}}{1-q}\right)q^{d_0}
 \left(\frac{d_0}{1-q}+\frac{q}{(1-q)^2}\right)\le\eta.
\end{equation}
Set $d_{k+1}=\lceil bd_k\rceil$. For each sufficiently large integer $N$, let
$K$ be maximal with $D'=\sum_{k=0}^{K-1}d_k\le\delta N/2$.
For each $k$ take the full ring of $2d_k$ numbers $z$ determined by
\begin{equation}\label{eq:ring}
 \prod_{z\text{ in ring }k}(x+zy)
 =x^{2d_k}+\epsilon_k\left(\frac{N}{2Cd_k}\right)^{d_k}y^{2d_k},
 \qquad \epsilon_k\in\{-1,1\}.
\end{equation}
Their squared moduli are $N/(2Cd_k)$. Add $N-2D'$ zero slopes.
Normalize every row $(1,z_i)$ to unit length. The matrix $A$ is a rank-two correlation
matrix: a reserve row exists and every ring has a nonzero slope. As each ring
has at least four points,
\begin{equation}\label{eq:slopes}
 \sum_i z_i=\sum_i z_i^2=0,
 \qquad \sum_i|z_i|^2=\frac{NK}{C},
 \qquad \sum_i(\Imm z_i)^2=\frac{NK}{2C}.
\end{equation}

Let $\gamma>0$ be the product of the first row coordinates. Write their product
polynomial as
\[
 P/\gamma=\sum_{j=0}^{D'}c_jx^{N-2j}y^{2j},\qquad
 R=\frac{\per A}{N!\gamma^2}
   =\sum_{j=0}^{D'}\frac{|c_j|^2}{\binom N{2j}}\ge1.
\]
Choose the signs in \eqref{eq:ring} independently and uniformly. For a subset
$S$ of ring indices, put $j(S)=\sum_{k\in S}d_k$. Its coefficient contribution
has magnitude
\[
 a_S=\left(\frac{N}{2C}\right)^{j(S)}
           \prod_{k\in S}d_k^{-d_k}.
\]
Distinct subsets have orthogonal sign products, including when their degree
sums coincide. Hence
\begin{equation}\label{eq:randomnorm}
 \mathbb ER=1+\sum_{S\ne\varnothing}
          \frac{a_S^2}{\binom N{2j(S)}}.
\end{equation}
By Lemma~\ref{lem:entropy}, $a_S\le[NE(b)/(2Cj)]^j$ for $j=j(S)$.
The elementary integral comparison
$\sum_{r=1}^h\log r\le\int_1^h\log t\,dt+\log h$ gives
$h!\le eh(h/e)^h$. For $2j\le\delta N$ it follows that
\[
 \binom N{2j}\ge\frac{[e(N-2j)/(2j)]^{2j}}{2ej}
 \ge\frac{[e(1-\delta)N/(2j)]^{2j}}{2ej}.
\]
Every summand in \eqref{eq:randomnorm} is therefore at most $2ejq^j$.
The subset generating function $Z(q)=\prod_k(1+q^{d_k})$ yields
\[
 \sum_Sj(S)q^{j(S)}
 =Z(q)\sum_k\frac{d_kq^{d_k}}{1+q^{d_k}}
 \le\exp\!\left(\sum_kq^{d_k}\right)\sum_kd_kq^{d_k}.
\]
The $d_k$ are distinct integers at least $d_0$, so bounding these sums by the
complete integer tails and using \eqref{eq:d0} proves $\mathbb ER\le1+\eta$.
At least one of the finitely many sign choices consequently satisfies
\begin{equation}\label{eq:Rsmall}
 R\le1+\eta.
\end{equation}
This is finite averaging, not an assertion about almost-sure limits.

For any complex vector $w$, define
\[
 G_w=\sum_i\overline{w_i}v_i\otimes\prod_{j\ne i}\ell_j.
\]
By \eqref{eq:fock}, $\|G_w\|^2=w^*C(A)w$. The coefficient of $x^{N-1}$ in
its second vector component is $\gamma\sum_i\overline{w_i}z_i$. Keeping only
this squared coefficient proves
\begin{equation}\label{eq:marked}
 w^*C(A)w\ge(N-1)!\gamma^2
                  \left|\sum_i\overline{w_i}z_i\right|^2.
\end{equation}
Use $w=z$ for the complex test. For the real test use $w_i=\Imm z_i$;
\eqref{eq:slopes} gives $\sum_iw_iz_i=i\|w\|^2$. Equations
\eqref{eq:Rsmall}--\eqref{eq:marked} give
\[
 \frac{\lambda_{\max}(C(A))}{\per A}\ge\frac{K}{C(1+\eta)},\qquad
 \frac{\lambda_{\max}(\Ree C(A))}{\per A}\ge\frac{K}{2C(1+\eta)}.
\]
The recurrence implies
$d_0b^k\le d_k\le(d_0+1/(b-1))b^k$, and maximality of $K$ therefore gives
$K=\log N/\log b+O(1)$, with all parameters fixed. The lower limits are at
least $1/[C(1+\eta)\log b]$ and half this. Let $C$ approach its allowed lower
endpoint, then let $\delta,\eta\downarrow0$ and $b\downarrow1$. Since
$E(b)\log b\to e$, they approach $1$ and $1/2$, matching \eqref{eq:upper}.
This proves the four limits in Theorem~\ref{thm:main}. For each desired accuracy,
parameters are fixed first and only then is $N$ taken to infinity; the lower
construction works for every sufficiently large $N$.

Positive diagonal congruence multiplies $C(A)$ and $\per A$ by the same
positive scalar. Thus correlation normalization loses no generality when
$\per A>0$. Finally $(A+\varepsilon I)/(1+\varepsilon)$ is positive-definite
and correlation whenever $A$ is correlation and $\varepsilon>0$. Continuity
at each finite $N$ retains any strictly smaller Rayleigh bound. The
positive-definite assertion does not retain rank two for $N>2$.

\section{An explicit dyadic family and exact certificates}
Fix $K\ge1$, let $M=2^K$, $n=2M$, and $c>4$. For $k=0,\ldots,K-1$, put
$d_k=2^k$, $r_k^2=(M-1)/(cd_k)$, and take the slopes
\[
 z_{k,j}=r_k\exp\left(\frac{(2j+1)\pi i}{2d_k}\right),
 \quad0\le j<2d_k,
\]
together with two zero slopes. Normalize the corresponding rows $(1,z)$.
Their product is
\[
 P=\gamma x^2\prod_k(x^{2d_k}+r_k^{2d_k}y^{2d_k})
  =\gamma\sum_{j=0}^{M-1}c_jx^{n-2j}y^{2j}.
\]
Binary expansion is unique. Lemma~\ref{lem:entropy} with $b=2$, for which
$E(2)=4$, gives
\[
 c_j\le(4/c)^j\binom{M-1}{j},\qquad
 1\le R=\sum_j\frac{c_j^2}{\binom{2M}{2j}}
     <\frac{c^2}{c^2-16}.
\]
The last bound uses Vandermonde's inequality
$\binom{2M}{2j}\ge\binom Mj^2$ and a geometric series. Root sums give
\[
 \sum|z_i|^2=\frac{2K(M-1)}c,\quad
 \sum z_i^2=-\frac{2(M-1)}c,\quad
 \sum(\Imm z_i)^2=\frac{(K+1)(M-1)}c.
\]
Applying \eqref{eq:marked} to $z$ and $\Imm z$ proves
\begin{equation}\label{eq:dyadic}
 \frac{\lambda_{\max}(C(A_n))}{\per A_n}
 >\frac{(c^2-16)K(M-1)}{c^3M},\quad
 \frac{\lambda_{\max}(\Ree C(A_n))}{\per A_n}
 >\frac{(c^2-16)(K+1)(M-1)}{2c^3M}.
\end{equation}
Taking $c=7$ gives the constants $33/343$ and $33/686$.

For reproducible finite values we give the full Rayleigh formulas. Put
$B=2(M-1)/c$, $c_M=0$, and let $s_2(j)$ be the number of ones in the binary
expansion of $j$. Define degree-$n-2$ polynomials
\[
 S_0=\gamma\sum_{j=0}^{M-1}B(K-s_2(j))c_jx^{n-2-2j}y^{2j},\quad
 S_1=\gamma\sum_{j=0}^{M-1}2(j+1)c_{j+1}x^{n-2-2j}y^{2j}.
\]
Terms with zero coefficients at the endpoint are omitted. Summing reciprocals
within each root ring gives
\[
 G_z=(-y,x)S_0,\qquad G_{\bar z}=(y,-x)S_1,
 \qquad G_{\Imm z}=\frac1{2i}(y,-x)(S_0+S_1).
\]
For example, on a ring of size $m$ and squared radius $a$,
$\sum\bar z/(x+zy)=-am,yx^{m-2}/(x^m+a^{m/2}y^m)$ and
$\sum|z|^2/(x+zy)=am,x^{m-1}/(x^m+a^{m/2}y^m)$.
The formula for $S_1$ follows by differentiating $P$ in $y$ and dividing by
$y$. Finally $\|xS\|^2+\|yS\|^2=n\|S\|^2$ for degree $n-2$.
These identities reduce both full Rayleigh quotients to rational arithmetic.
The two independent supplied programs use, respectively, binary coefficients
and a marked-product recurrence. Exact integer calculations certify
\[
 \begin{array}{c|c|c}
 (K,c,n)&\text{complex direction }z&\text{real direction }\Imm z\\\hline
 (8,7,512)&1137/1000<Q<1138/1000&639/1000<Q<640/1000\\
 (10,5,2048)&1991/1000<Q<1992/1000&1094/1000<Q<1095/1000.
 \end{array}
\]
These finite examples illustrate the formulas, not the proof of the asymptotic
limits. The actual normalized matrix entries are algebraic; the Rayleigh
quotients are rational because the common row-normalization factor cancels.

\section{The sharp ramp constant}
Let $w_i=2i-n-1$ and define $M=C(A)/\per A$. For $1\le j<n$, let
$b_j(i)=\one_{\{i>j\}}$ and $u_j=b_j-(n-j)\one/n$. Lemma~\ref{lem:indicator}
and $M\one=\one$ give
\[
 u_j^*Mu_j\le\frac{j(n-j)}n,
 \qquad w=2\sum_{j=1}^{n-1}u_j,
 \qquad\|w\|^2=\frac{n(n^2-1)}3.
\]
The triangle inequality in the $M^{1/2}$ seminorm proves the all-rank bound
\begin{equation}\label{eq:ramp}
 \frac{w^*C(A)w}{\per A\,\|w\|^2}
 \le B_n:=\frac{12\bigl(\sum_{j=1}^{n-1}\sqrt{j(n-j)}\bigr)^2}
                 {n^2(n^2-1)}
 \longrightarrow\frac{3\pi^2}{16}.
\end{equation}
The limit follows from the Riemann sum for
$\int_0^1\sqrt{t(1-t)}\,dt=\pi/8$.

We include the lower construction, rather than infer sharpness from numerics.
Choose unit row configurations whose empirical projective measures approach
Haar measure on $\CP^1$, for example uniform grids in squared first-coordinate
modulus and relative phase. Let $p_m$ be any global projective maximizer of the
product of the $m$ forms. No form vanishes there. Apply a unitary variable change
sending $p_m$ to $(1,0)$ and write the resulting ratios as $z_i=b_i/a_i$.
The rotated empirical measures still approach Haar measure: every sequence of
unitaries has a convergent subsequence, and Haar measure is invariant under
its limit. The ratio map is continuous outside one Haar-null point.
Thus the empirical real parts approach the law
\[
 f_X(x)=\frac1{2(1+x^2)^{3/2}},\qquad
 F_X(x)=\frac12\left(1+\frac{x}{\sqrt{1+x^2}}\right).
\]
For independent $X,Y$ with this law,
\[
 \mathbb E|X-Y|=2\int F_X(x)(1-F_X(x))\,dx=\frac\pi2.
\]
Weak convergence of product empirical measures applied first to the bounded
continuous function $\min(|x-y|,T)$, and then monotone convergence as
$T\to\infty$, give
\begin{equation}\label{eq:pairs}
 \liminf_m\frac1{m^2}\sum_{i<j}|\Ree z_i-\Ree z_j|\ge\frac\pi4.
\end{equation}
This does not assume convergence of unbounded first moments.

Append the unit row form $p_m^*u$. It has modulus at most one on the sphere,
with equality only at the projective point $p_m$. The new product consequently
has a unique maximum there. In the rotated coordinates this adds one zero
ratio, and \eqref{eq:pairs} persists with $n=m+1$. Order these $n$ rows by
increasing $\Ree z_i$. Stationarity of the maximum gives $\sum_i z_i=0$.

For each fixed augmented base repeat every row $L$ times consecutively, let
$N=nL$, and let $U_L$ be the isometry taking a base coordinate to the normalized
indicator of its cluster. Write $P$ for the base product. The sphere identity
\[
 \langle f,g\rangle_B=(d+1)!\int_{\CP^1}f\bar g\,d\mu
 \quad(\deg f=\deg g=d)
\]
follows by phase integration and the beta integral for monomials. It gives the
exact compressed normalized cofactor matrix
\[
 K_L:=\frac{U_L^*C(A_L)U_L}{\per A_L},\qquad
 (K_L)_{ij}=\frac{L}{nL+1}\mathbb E_{\mu_L}
       \frac{\langle v_i,v_j\rangle}{\ell_i\overline{\ell_j}},
 \quad d\mu_L\ \propto\ |P|^{2L}d\mu.
\]
The unique maximum implies concentration at $(1,0)$. The kernels' poles are
harmless: outside a fixed neighborhood of the maximum,
\[
 |P|^{2L}\frac1{|\ell_i\ell_j|}
 =|P|^{2L-2}|P/\ell_i||P/\ell_j|.
\]
The quotients are bounded polynomials on the sphere, and the maximum of
$|P|$ on that complement is strictly smaller than its value on a sufficiently
small neighborhood of the peak. Comparing with that neighborhood in the
normalizing integral proves exponential decay of the displayed tail. Locally
the kernels are continuous. Hence
\[
 K_L\longrightarrow\frac{J+zz^*}{n}.
\]
All cofactor blocks between row clusters are constant, so $C(A_L)$ annihilates
$\operatorname{ran}(U_L)^\perp$. The order-$N$ ramp compresses exactly to
$L^{3/2}$ times the order-$n$ ramp. Its Rayleigh quotient therefore tends to
\[
 \frac{3|\sum_i(2i-n-1)z_i|^2}{n^4}.
\]
The real part of the numerator's unsquared sum equals
$\sum_{i<j}|\Ree z_i-\Ree z_j|$. By \eqref{eq:pairs}, a diagonal sequence of
finite base sizes and sufficiently large finite $L$ attains the lower limit
$3\pi^2/16$. Together with \eqref{eq:ramp}, this proves
\[
 \limsup_{N\to\infty}\sup_{\rank A=2}
 \frac{w_N^*C(A)w_N}{\per A\,\|w_N\|^2}=\frac{3\pi^2}{16}.
\]
The same limsup holds without the rank restriction and within positive-definite
correlation matrices by continuity. No dimension-by-dimension limit of these
ramp suprema is asserted.

\section{Rank-two endpoint interpretation}
For binary row forms $\ell_i=a_ix+b_iy$ put
\[
 S=\sum_{i<j}(a_ib_j-b_ia_j)\prod_{r\notin\{i,j\}}\ell_r.
\]
For the real ramp weights,
\[
 G_w=\sum_{i<j}(v_j\ell_i-v_i\ell_j)
                \prod_{r\notin\{i,j\}}\ell_r=(-y,x)S.
\]
Thus the Fock identity gives
$n\|S\|^2=w^*C(A)w$. Substitution in \eqref{eq:ramp} yields
\[
 \frac{\|S\|^2}{\binom n2\per A}
 \le D_n:=\frac{8\bigl(\sum_{j=1}^{n-1}\sqrt{j(n-j)}\bigr)^2}
                 {n^3(n-1)}\longrightarrow\frac{\pi^2}{8},
\]
and the preceding rank-two construction supplies the matching lower limsup.

For clarity, if $P_q(A)=\sum_{\sigma\in S_n}q^{\operatorname{inv}(\sigma)}
\prod_i a_{i,\sigma(i)}$, direct expansion also gives
\[
 2P'_1(A)=\binom n2\per A-\|S\|^2.
\]
Indeed the product of two $2$ by $2$ row determinants is the corresponding
$2$ by $2$ Gram determinant. Expanding $\|S\|^2$, a permutation receives a
positive contribution from each non-inversion and a negative contribution
from each inversion, for total coefficient $\binom n2-2\operatorname{inv}(\sigma)$.
Consequently $\pi^2/8$ is the optimal asymptotic constant for this normalized
rank-two endpoint ratio. The identity linking $S$ to the ramp cofactor form
has only been used in rank two; no all-rank endpoint bound with this constant
is claimed. The endpoint counterexample mechanism itself was already
established in the preceding research; the matching sharp upper constant is
the issue treated here.

\section{What remains unresolved}
Neither a large noncentral cofactor eigenvalue nor its real Rayleigh excess
is automatically a normalized character or subgroup projector. Ordinary
rank-two immanants lie in two-row sectors where classical dominance results
already apply. Marcus's rank-two large-order regime is independently excluded
by a factorial/Jensen comparison; the precise inequality and finite thresholds
are recorded in the stage summary. Higher-rank and genuine subgroup couplings
have not been resolved. Numerical searches are reported only as searches.

An unrestricted higher-compound binary-vector formula was separately tested
under an explicit definition and has an exact internal counterexample. The
full text and qualifications of the potentially related statement in
\cite{pate08} were not obtained. We make no assertion that its theorem is
false, and do not use that unverified statement anywhere in this paper.

\clearpage

\section*{Source edition and reproduction}
This signed article edition retains the mathematical argument of the
previously archived source: \url{https://doi.org/10.5281/zenodo.23248496}.
The original source, exact inputs, checking programs and prior review
records remain available there. The companion README identifies the
frozen Git version and the scope of each auxiliary computation.
In a writable copy of that archived source directory, run:
\begin{verbatim}
python3 verify_dyadic_cofactor_family.py --K 8 --C 7
python3 verify_dyadic_cofactor_family.py --K 10 --C 5
python3 verify_dyadic_independently.py
python3 verify_stage_bounds.py
\end{verbatim}
The computational checks do not replace the analytic proof.

\section*{Funding and research methods}
This research was conducted independently and received no external
funding. OpenAI AI assistants assisted mathematical research, verification
code and internal review. Mathematical claims are supported by the
written proof and the specified public evidence; AI output is not itself
proof evidence. Internal AI-assisted review is not external human peer
review.
\section*{Declaration of generative AI and AI-assisted technologies
in the writing process}
OpenAI AI assistants assisted drafting, organization and LaTeX editing.
The author takes responsibility for the manuscript and its descriptions
of the proof and verification evidence.

\begin{thebibliography}{9}
\bibitem{drury} S. W. Drury, A counterexample to a question of Bapat \& Sunder,
\emph{Math. Inequal. Appl.} 21 (2018), 517--520.
\href{https://doi.org/10.7153/mia-2018-21-37}{doi:10.7153/mia-2018-21-37}.
\bibitem{pate96} T. H. Pate, Group algebras, monotonicity, and the Lieb permanent
inequality, \emph{Linear Multilinear Algebra} 40 (1996), 207--220.
\href{https://doi.org/10.1080/03081089608818438}{doi:10.1080/03081089608818438}.
\bibitem{pate08} T. H. Pate, On permanental compounds,
\emph{Linear Algebra Appl.} 429 (2008), 1093--1101.
\href{https://doi.org/10.1016/j.laa.2007.05.019}{doi:10.1016/j.laa.2007.05.019}.
\bibitem{neuberger} J. W. Neuberger, Norm of symmetric product compared with norm
of tensor product, \emph{Linear Multilinear Algebra} 2 (1974), 115--121.
\href{https://doi.org/10.1080/03081087408817047}{doi:10.1080/03081087408817047}.
\bibitem{pioge} L. Pioge, K. K. Pietrasz, B. Seron, L. Novo, N. J. Cerf,
A logical implication between two conjectures on matrix permanents,
\emph{Linear Algebra Appl.} 725 (2025), 309--318.
\href{https://doi.org/10.1016/j.laa.2025.07.011}{doi:10.1016/j.laa.2025.07.011};
\href{https://arxiv.org/abs/2508.00111v1}{arXiv:2508.00111v1}.
\bibitem{wanless} I. M. Wanless, Lieb's permanental dominance conjecture,
in \emph{The Physics and Mathematics of Elliott Lieb}, Vol. II,
EMS Press (2022), 501--516.
\href{https://doi.org/10.4171/90-2/48}{doi:10.4171/90-2/48};
\href{https://arxiv.org/abs/2202.01867v1}{arXiv:2202.01867v1}.
\end{thebibliography}
\end{document}
